\section{March 09}

\subsection{Introduction(continue)}

    \Yang{This subsection has many typos and errors. I will try to fix them when I have time and understand the content better.}

    Recall: Let \(Z = Z(s) \subset \bbP^4, s \in H^0(\calO_{\bbP^4}(5))\) be a generic nember.
    Let \(n_1 = \#\{\text{lines in } Z\}\).

    \begin{exercise}
        Show that \(\calN_{L/Z} \cong \calO_{\bbP^1}(-1) \oplus \calO_{\bbP^1}(-1)\).
    \end{exercise}

    \(n_1\) can be computed by 
    Schubert calculus
    \begin{align*}
        n_1 &= \int_{\Gr(2,5)} c_6(\Sym^5 \calS^\vee) \\
            &= \int_{\Gr(2,5)} \left( 600 \sigma^4_\square \sigma^2_{\square\square} + 1450 \sigma^2_{\square} \sigma^2_{\square\square} + 225 \sigma_{\square\square}^3 \right) \\
            &= 2875.
    \end{align*}
    or by equivariant localization
    \begin{align*}
        n_1 &= \sum_{I=\{i_1,i_2\} \subset \{1,2,3,4,5\}, |I|=2} \frac{\prod_{a = 0}^5 (a\lambda_{i_1} + (5-a)\lambda_{i_2})}{ \prod_{i,j \notin I} (t_i - t_j)} \\
            &= 2875.
    \end{align*}

    And the number \(n_2\) of conics in \(Z\) can be computed by
    \begin{align*}
        n_2 &= \int_{\bbP(\Sym^2 \calS^\vee) \to \Gr(3,5)} c_{11}(\calE) \\
            &= 609250.
    \end{align*}
    where \(\calE = (\pi^* \Sym^5 \calS^\vee)/ (\calU \ten \pi^*\Sym^3 \calS^\vee)\).

    It can be computed by Maple.

    Let 
    \[ Z_0 = \{x_0^5 + x_1^5 + x_2^5 + x_3^5 + x_4^5 = 0\} \subset \bbP^4. \]
    Let \(\zeta\) be a primitive 5-th root of unity.
    Let $H_{ij\zeta} := \{x_i - \zeta x_j = 0\} \subset \bbP^4$ for any \(i \neq j\).
    Let $C_{ij}$ be the plane Fermat conic defined by 

    e.g. \(X = (u,-\zeta^k u, av, bv,cv), a^5 + b^5 + c^5 = 0\).
    Then \(C_{ij} \subset Z_0\) and \(C_{ij} \subset H_{ij\zeta}\) for some \(\zeta\).
    We get \(1\)-parameter family of lines in \(Z_0\).

    There are \(5\) choice of \(k\), \(C_5^2\) choice of \(\{i,j\}\), so we get \(5 \cdot C_5^2 = 50\) such families.
    These cone meet in \(5 \cdot (1/2) C_4^2 5^3 = 375\) isolated lines of the form $(u, -\zeta^k u, v, -\zeta^l v, 0)$.
    Each cone contains 15 isolated lines and meets 15 other cones in these lines.

    Theree are no other intersection points.

    \begin{exercise}
        For such a line \(L\), show that \(\calN_{L/Z_0} \cong \calO_{\bbP^1}(-3) \ten \calO_{\bbP^1}(1)\).
    \end{exercise}

    One can show that a line on \(Z_0\) is contained in one of these families.
    \(\implies\) The space of lines on \(Z_0\) has codimension \(1\).

    \begin{question}
        How do we count lines on \(Z_0\)?
    \end{question}

    Solution: under a sufficiently general deformation of $Z_0$, 
    \begin{itemize}
        \item each isolated splits into 5 lines.
        \item each cone splits into 20 lines.
    \end{itemize}
    Hence we get \(375 \cdot 5 + 50 \cdot 20 = 2875\).

    \paragraph{The 4th method to get \(n_1=2875\) from physics.(Candelas-de la Ossa-Green-Parkes, 1991)}

    As before,  \(Z \subset \bbP^4\) is a smooth variety over \(\bbC\) with \(K_Z = 0, h^{1,0} = 0\).
    i.e. \(Z\) is a Calabi-Yau 3-fold.
    Moreover, we have \(h^{1,1} = 1\) and \(h^{2,1} = 101\).
    Physics propose:
    there exists \(f:\calZ^\vee \to \bbP^1\) family of Calabi-Yau 3-folds with \(h^{1,1}(\calZ_t^\vee) = 101\) and \(h^{2,1}(\calZ_t^\vee) = 1\) for all \(t \in \bbP^1\).

    There is an associated polarized VHS $(R^3f_* \bbZ_{\calZ^\vee}, \calF^\bullet, Q)$ over \(\bbP^1\) with a canonical flat connection \(\nabla\).
    Let \(T_t \in \Aut(H^3(\calZ_t^\vee, \bbZ))\) be the monodromy of \(\nabla\) around \(t\).
    Then there is a unique point \(0 \in \bbP^1\) such that \((T_0 - 1)^4 = 0\) but \((T_0 - 1)^3 \neq 0\).
    There is a unique flat section \(s = (s_0,\ldots,s_4)\) near \(0\) such that 
    \[ s_0 \in 1+z \bbC[[z]], s_1 \in s_0\log z + z \bbC[[z]], s_2 \in s_0 \frac{\log^2 z}{2} + s_1 \log z + z \bbC[[z]]. \]
    Equivalently, \(s_i(z)\) are solution to the ODE
    \[ \left( \theta^4 - 5 z \prod_{k=1}^4 ( \theta + k ) \right) s = 0,\quad \theta = z \frac{d}{dz}. \]

    Define a new coordinate on \(\bbP^1\) near \(0\) by \(q = q(z) = \exp(s_1(z)/s_0(z)) \in z\bbC[[z]]\).
    Invert this series to get \(z = z(q) \in q \bbC[[q]]\).
    There is a rational function on $\bbP^1$, the normalized Yukawa coupling, given as
    \[ C = 5 \theta_q^2\left( \frac{s_2}{s_0}(z(q)) \right) \in 5 + q \bbC[[q]], \quad \theta_q = q \frac{d}{dq}. \]
    Consider the expansion of \(C\) in a Lambert series
    \[ C = \sum_{d \geq 0} \tilde{n}_d d^3 \frac{q^d}{1-q^d}. \]
    Computing the series, one can get \(\tilde{n}_1 = 2875\).

    Using equivariant localization, Ellingson and Shornme computed \(n_3\) and reproduced this prediction of physics.

    (1994) Kontserich proposed to study the moduli space $\calM_{0,0}(Z,\beta)$ of stable maps \(f:\bbP^1 \to Z\) such that \(f(\bbP^1) = \beta\) for some effective \(\beta \in H_2(Z,\bbZ)\).
    And applied equivariant localization on $\overline{\calM_{0,0}}(\bbP^4, \beta)$ and restricted to $\overline{\calM}_{0,0}(Z,\beta)$.
    Together with Manin, they set up a system of axioms defining \emph{Gromov-Witten invariants}.

    (1996-98) the procedure of Canclelas et al. was proven by Girental and Licm-Liu-Yau.

    What are the general lessons from these examples:
    \begin{itemize}
        \item Find a moduli space \(\calM\) of curves (of stable maps from curves).
        \item Find a nice space \(\calX\) (smooth scheme of DM stack) s.t. $\calM$ is a local complete intersection in $\calX$;for example, $M=Z(s), s \in H^0(\calU,\calE), \calU \subset \calX$ open.
        \item \(\calE\) is obtained from the deformation theory of objects of \(\calM\).
        \item the invariant is obtained as 
            \[ e(\calE) \cap [\calX] = i_*(s_*[C_\calX \calM]). \]
        \item if \(\calM\) admits a \(\bbC^*\) action, this class can be computed by equivariant localization.
        \item The naive $[\calX]$ can have components of different dimensions. Need a generalization of the notion of a fundamental class.
    \end{itemize}


\subsection{Review of moduli of curves}

    A (complex) curve is a variety $C \to \Spec \bbC$ of dimension $1$, i.e. a reduced, separated scheme of finite type over $\Spec \bbC$ such that all irreducible components have dimension $1$.
    If $C$ is projective, its arithmetic genus of $C$ is $p_a(C) = 1- \chi(C,\calO_C)$.

    \begin{definition} 
        A closed point $q \in C$ is a node if the following equivalent conditions hold:
        \begin{itemize}
            \item there exists a neighborhood of $q \in C(\bbC)^\an$ which is complex-analytically isomorphic to a neighborhood of $0$ in $\{(x,y) \in \bbC^n | xy = 0\}$;
            \item or the completion of $\calO_{C,q}$ at $\frakm_q$ satisfies $\hat{\calO}_{C,q} \cong \bbC[[x,y]]/(xy)$.
            \item or the normal cone $N_q C$ is a union of two reduced lines.
        \end{itemize}
        The curve $C$ is \emph{nodal} if every closed points $q \in C$ is either smooth points or a node.
    \end{definition}
   
    \begin{definition}
        A nodal curve $C$ together with an ordered collection of $n$-points $p_1,\cdots,p_n \in C$ is called an \emph{$n$-pointed curve}, $p_i$ is called \emph{marked points}.
    \end{definition}

    \begin{definition}
        A connected, nodal, projective curve $C$ is \emph{stable} if $\Aut(C)$ is finite.
        A connected nodal projective $n$-pointed curve $(C,p_1,\cdots,p_n)$ is \emph{stable} if $\Aut(C,p_1,\cdots,p_n) = \{f \in \Aut(C) \mid f(p_i) = p_i \text{ for all } i\}$ is finite.
    \end{definition}
    
    \begin{proposition}
        Let $(C,p_1,\cdots,p_n)$ be an $n$-pointed nodal projective connected curve, and $\tilde{C} = \sqcup \tilde{C}_i$ be the normalization of $C$.
        Then $(C,p_1,\cdots,p_n)$ is stable if and only if for each $i$, one of the following holds:
        \begin{enumerate}
            \item $g(\tilde{C}_i) = 0$ and $\tilde{C}_i$ has at least $3$ special points (i.e. marked points or preimages of nodes);
            \item $g(\tilde{C}_i) = 1$ and $\tilde{C}_i$ has at least $1$ special point;
            \item $g(\tilde{C}_i) \ge 2$.
        \end{enumerate}
    \end{proposition}

    \begin{definition}
        An $n$-pointed family of smooth (resp. stable) curves over a scheme $S$ is a tuple $(\phi:\calC \to S, \sigma_1,\cdots,\sigma_n:S \to \calC)$ where
        \begin{enumerate}
            \item $\phi$ is a smooth (resp. flat) proper surjective finitely presented morphism of schemes s.t. every geometric fiber is a smooth (resp. stable) $n$-pointed projective connected curve of arithmetic genus $g$.
            (If $C_s$ is only nodal, we call $\phi:\calC \to S$ prestable.)
            \item $\sigma_i$ are pairwise disjoint sections of $\phi$ with image in the smooth locus of $\phi$.
        \end{enumerate}
    \end{definition}

    Let $\alpha:T \to S$ be a morphism of schemes, $(\calC/T, \sigma_{1,T}, \cdots, \sigma_{n,T})$ be the pullback of $(\calC/S, \sigma_1,\cdots,\sigma_n)$ along $\alpha$.

    \begin{definition}
        Two $(\calC,\sigma_1,\cdots,\sigma_n)/S$ and $(\calC',\sigma_1',\cdots,\sigma_n')/S$ are isomorphic if there exists an $S$-isomorphism $\calC \cong \calC'$ sending $\sigma_i$ to $\sigma_i'$ for all $i$, i.e. there exists an isomorphism $\calC \cong \calC'$ such that the following diagram commutes:
        \[\begin{tikzcd}
            \calC \ar[rr, "\cong"] \ar[dr, "\phi"] & & \calC' \ar[dl, "\phi'"'] \\
            & S \ar[ur, "\sigma_i'"', bend right] \ar[ul, "\sigma_i", bend left] &
        \end{tikzcd}\]
    \end{definition}
    
    \begin{definition}
        Define the category $\overline{\calM}_{g,n}$ over $\Sch$ as follows:
        \begin{itemize}
            \item \(\Obj(\overline{\calM}_{g,n}) = \{(S, (\calC/S, \sigma_1,\cdots,\sigma_n))\}\), where $S$ is a scheme and $(\calC/S, \sigma_1,\cdots,\sigma_n)$ is an $n$-pointed family of genus $g$ stable curves over $S$.
            \item \(\Hom_{\overline{\calM}_{g,n}}((S, (\calC/S, \sigma_1,\cdots,\sigma_n)), (S', (\calC'/S', \sigma_1',\cdots,\sigma_n')))\) is the set of morphisms $(f,\alpha)$ with diagram 
            \[\begin{tikzcd}
                \calC \ar[r, "f"] \ar[d, "\phi"']  & \calC' \ar[d, "\phi'"] \\
                S \ar[r, "\alpha"] & S'
            \end{tikzcd}\]
            being Cartesian, and $f \circ \sigma_i = \sigma_i' \circ \alpha$ for all $i$.
            \Yang{In the blackboard, the arrows are from $(S',\calC')$ to $(S,\calC)$, one need to be careful about the direction.}
        \end{itemize}
        Similarly, we can define $\calM_{g,n}$ (resp. $\calm_{g,n}$) of smooth (resp. prestable) curves.
    \end{definition}

    \begin{definition}
        Let $\calC$ be a site.
        We say that $\calF \in \Hom(\calC^\op, \Grpd)$ is a \emph{stack} if for any covering $U_i \to U$ of any $U \in p(\calF)$, 
        \begin{enumerate}
            \item given $x_i \in \calF(U_i)$, $\phi_{ij}:x_i|_{U_ij} \to x_j|_{U_{ij}}$ such that $\phi_{ij}|_{U_{ijk}} \circ \phi_{jk}|_{U_{ijk}} = \phi_{ik}|_{U_{ijk}}$,
                there exist $x \in \calF(U)$ and isomorphisms $\phi_i:x|_{U_i} \cong x_i$ such that $\phi_{ji} \circ \phi_i|_{U_{ij}} = \phi_j|_{U_ij}, \forall i,j \in I$;
            \item for all $x,x' \in \calF(U)$, the presheaf 
                \[ Isom_{U}(x,x'): (U_i \to U) \mapsto \{ \text{isomorphisms } x|_{U_i} \to x'|_{U_i} \} \]
                is a sheaf on $\calC/U$,
                or equivalently, 
                \begin{enumerate}
                    \item $\forall x,x' \in \calF(U), \phi_i: x|_{U_i} \cong x'|_{U_i}$ such that $\phi_i|_{U_{ij}} = \phi_j|_{U_{ij}}$, there exists a unique $\phi:x \cong x'$ such that $\phi|_{U_i} = \phi_i$ for all $i$;
                    \item given $x,x' \in \calF(U), \phi:x \to x', \psi:x \to x'$ such that $\phi|_{U_i} = \psi|_{U_i}$, then $\phi = \psi$.
                \end{enumerate}
        \end{enumerate}
    \end{definition}

    Let 
    \[ \begin{tikzcd}
        & \calF_1 \ar[d] \\
        \calF_2 \ar[r] & \calF_3
    \end{tikzcd} \]
    be a diagram of stacks over $\calC$.
    The fiber product $\calF_1 \times_{\calF_3} \calF_2$ is also a stack over $\calC$.



